\begin{proposition}[Semi-log-canonical fibre supports]\hypertarget{tag-02WX}{}\label{night:P4.25}\label{stmt:P4.25}
\StatementRecord{night:P4.25}{proved}{02WX}{\SeriesRef{night:P4.16}, \SeriesRef{night:P4.20}, \SeriesRef{night:P4.6}, \SeriesRef{night:P4.8}}

Let $f:Y\to B$ be a proper surjective holomorphic map with connected
fibres from a connected smooth compact complex threefold to a smooth
compact curve. Assume $b_1(Y)=b_2(Y)=0$ and that the smooth general
fibre is a complex two-torus. Every semi-log-canonical reduced fibre
support of $f$ has normal crossings.

For an irreducible reduced support $S$, let $\nu:N\to S$ be its
normalization, let $\pi:T\to N$ be a resolution, and put
$\rho=\nu\circ\pi$. For local residue generators $\Omega$ of $\omega_S$,
define the integral divisor $D=-\operatorname{div}(\rho^*\Omega)$, so that
$\omega_T(D)\simeq\rho^*\omega_S$ with this identification.
A reduced normal-crossing divisor $D$ gives normal crossings on $S$
and a smooth normalization $N$. Proposition~\ref{night:P4.6} applies
to each such support. Propositions~\ref{night:P4.7}--\ref{night:P4.8}
apply to fibrations whose reduced supports are all semi-log-canonical.
\end{proposition}
\begin{proof}

\begin{enumerate}[label={$\langle1\rangle$\arabic*.},leftmargin=*]
\item At a generic point of the conductor, an exponent one on every normalization
branch forces the transverse plane curve to be an
ordinary node. Indeed the conductor exponent on
branch $i$ is
$2\delta_i+\sum_{j\ne i}I_{ij}$; positivity of
intersection multiplicities leaves exactly two
smooth transverse branches. The support $S$ is a
Cartier hypersurface, hence is Cohen--Macaulay and
has invertible dualizing sheaf. On its normalization $N$, the pair consisting of $N$
and its reduced conductor is log canonical.
To check this last assertion, pull its canonical
divisor to $T$. The resulting boundary is exactly
$D$, by the displayed isomorphism and the
divisorial orders of a local generator. A smooth
surface with reduced normal-crossing boundary is
log canonical: blowing up a point gives new
boundary coefficient $\mu-1\le1$, and this
bound persists under successive point blowups.
Discrepancy comparison therefore gives the
assertion on $N$. Definition~2.3 and Lemma~2.2 of
\cite[p.~189]{LiuRollenske2012SLC} give semi-log-canonicity in
the complex analytic category.

\item The analytic classification is
\cite[Theorem~1.1, p.~186; Table~1, p.~187;
Propositions~3.1 and~4.1, pp.~192--200]{LiuRollenske2012SLC}.
An isomorphism of the singular hypersurface germs lifts to an ambient
biholomorphism: lift its three coordinate classes to convergent functions;
the induced cotangent isomorphism makes their linear part invertible.
Thus, in the displayed coordinates, the full local fibre function is
$U\prod_i g_i^{m_i}$, with $U$ a unit and the original component
multiplicities $m_i$. The normal singular germs in the table are isolated.
Each is irreducible, since distinct surface branches would meet along
a singular curve. Write its full function as $Ug^m$, choose a holomorphic
$m$th root of $U$, and set $h=U^{1/m}g$, so that $f=h^m$. Any such germ is excluded
in our family: a small generic linear
perturbation $h-c\cdot z$ has Morse critical points.
Their number is the positive multiplicity of the gradient map
\cite[Appendix~B, Lemmas~B.1--B.3 and the positivity proof,
pp.~111--114]{Milnor1968}. It is also the colength of the gradient ideal:
the map $\mathbb C\{u\}\to\mathbb C\{z\}$,
$u_i\mapsto h_{z_i}$, is finite; its three parameters form a regular
sequence, so its target is a free module of that colength. A regular
nearby fibre consists of that many reduced points, each of local degree
one. The boundary degree is unchanged under the small perturbation. A constant perturbation
puts their critical values away from zero.
The original zero divisor is smooth near
a small transverse boundary sphere, so it
retains that boundary property under these
perturbations. The proper boundary
transport in the proof of Proposition~\ref{night:P4.20} carries a
vanishing sphere of square $-2$ into a
nearby torus fibre, which is impossible.

\item The ordinary double point and $xyz=0$ have normal crossings.
The change $z\mapsto-z$ identifies the pinch point $x^2+y^2z=0$
with the case $n=2$ of Proposition~\ref{night:P4.16}.
For a weighted homogeneous full function $G$ of positive weights
$(w_1,w_2,w_3)$ and degree $W$, the unit is absorbed by the
biholomorphism
\[
 (z_1,z_2,z_3)\longmapsto
 (U(z)^{w_1/W}z_1,U(z)^{w_2/W}z_2,U(z)^{w_3/W}z_3),
\]
using one local logarithm of $U$; it satisfies $G\circ\Psi=UG$.
The weights used below are positive. Scaling the coordinates by
$\lambda^{w_i}$ places every compact cycle in $G^{-1}(1)$ in
$G^{-1}(\lambda^W)$ inside an arbitrarily small neighborhood.
For a deformation $G_s$ of an irreducible reduced equation with full
multiplicity $m$, use the full map $G_s^m$. At nonzero values its
levels are the disjoint union of the levels $G_s=\zeta t^{1/m}$,
where $\zeta^m=1$, and
$d(G_s^m)=mG_s^{m-1}dG_s$. Thus the boundary transport on these
levels carries the original full multiplicity.
The remaining double-point equations are
\[
 x^2+y^2z^2,\qquad
 x^2+y^2(z^2+y^{q-2}),\quad q\ge3.
\]
For the first equation the two normalization
branches are the smooth graphs $x=\pm i yz$.
Their tangent planes coincide at the origin,
whereas they meet transversely off that point
along the two axes. Put $U_1=x-iyz$ and $U_2=x+iyz$, and let their full fibre
multiplicities be $a,b>0$. For $d=\gcd(a,b)$ write
$a'=a/d$, $b'=b/d$. On the component $U_1^{a'}U_2^{b'}=1$
of $U_1^aU_2^b=1$, set
\[
 U_1=t^{b'},\qquad U_2=t^{-a'},\qquad
 yz=\frac{t^{-a'}-t^{b'}}{2i}.
\]
The inverse parameter is $t=U_1^rU_2^s$ for integers
$rb'-sa'=1$. The last equation has $a'+b'\ge2$ simple nonzero
zeros in its right-hand side. Over an arc joining two of them,
the circles $|y|=|z|$ collapse at the endpoints and form a smooth
embedded sphere. The circle action $(y,z)\mapsto(e^{i\theta}y,
e^{-i\theta}z)$ has moment map proportional to $|y|^2-|z|^2$,
so the sphere is Lagrangian. Its oriented normal Euler number is
$-\chi(S^2)=-2$: the order $(e_1,e_2,Je_1,Je_2)$ differs by one
transposition from the complex orientation $(e_1,Je_1,e_2,Je_2)$.
Weights $(2,1,1)$ place this sphere in the full nearby fibre with
the original multiplicities. Every sphere in a complex two-torus
lifts to $\mathbb C^2$ and contracts there, so its homology class,
and hence its self-intersection, is zero. This excludes the germ.

\item For finite $q$, the equation is analytically irreducible:
its weighted homogeneous polynomial of weights $(q,2,q-2)$ is
irreducible, and lowest weighted terms of an analytic factorization
would factor that polynomial. Thus its full multiplicity is a common
integer $m$. After the unit map above, put
\[
 G_s=x^2+y^2(z^2+y^{q-2}+s).
\]
For $s\ne0$ small, at each point
$(0,0,z_0)$ with $z_0^2=-s$ the function
$Z=z^2+y^{q-2}+s$ is a local coordinate:
its $z$ derivative is $2z_0\ne0$.
The reduced equation in these coordinates is
$x^2+y^2Z$, a pinch point. At the original
boundary double axis $x=y=0$, $z\ne0$,
the transverse Hessian determinant is
$4(z^2+s)$, which stays nonzero on the
compact boundary when $s$ is small.
All boundary smooth points remain smooth.
Thus the zero divisor retains the same
boundary smooth and double strata.
The spherical class of square $-4$
from the pinch-point nearby fibre transports,
by the boundary transport in Proposition~\ref{night:P4.20}, to
the original nearby torus fibre.
This excludes every finite $q$.

\item The remaining triple-point equations are
\[
 xyz+x^p,\quad p\ge3,\qquad
 xyz+x^p+y^q,\quad q\ge p\ge3.
\]
For the first put
$G_s=xyz+x^p+s x^2
=x(yz+x^{p-1}+s x)$.
When $s\ne0$, its two local normalization
branches are smooth: the derivative in $x$
of the second factor at the origin is $s$.
Their tangent planes coincide. The two branches of $uv=0$ have
distinct tangent planes, and an ambient biholomorphism has an invertible
derivative preserving that distinction. Hence these two tangent branches
cannot be a normal crossing. Along
the boundary $y$ axis the transverse
Hessian determinant is $-y^2$, and along
the boundary $z$ axis it is $-z^2$.
Both axes persist as ordinary-double
strata. For full component multiplicities $a,b$, weights $(2,p-1,p-1)$
give degree $2a+2(p-1)b$ and absorb the original unit. Use the deformation
$x^a(yz+x^{p-1}+sx)^b$. At its boundary double strata, free normal
coordinates $u,v$ give $u^av^b$, which is submersive at nonzero
values; thus the same boundary transport applies with these multiplicities.
For fixed nonzero $s$, the coordinate $X=x(s+x^{p-2})$ is
biholomorphic near the origin. The full deformed function becomes
a unit times $X^a(X+yz)^b$. The coordinates
\[
 x'=X+yz/2,\qquad y'=y,\qquad z'=z/(2i)
\]
satisfy $x'-iy'z'=X$ and $x'+iy'z'=X+yz$.
After absorbing the unit with weights $(2,1,1)$, the preceding
two-branch case supplies a sphere of square $-2$. Boundary transport
therefore gives the contradiction in the original full fibre.

\item For the second equation, positive weights
$(q,p,pq-p-q)$ and the primitive degree-one expression in $z$
show analytic irreducibility by the same initial-form argument.
It therefore has one common full multiplicity. Put
\[
 G_s=xyz+x^p+y^q+s x^2+s y^3,
 \qquad 0<|s|\ll1.
\]
Its transverse Hessian along the boundary
$z$ axis has determinant $-z^2$,
so its sole boundary double stratum
persists. We check the smooth normalization
of its germ at the origin explicitly.
Since $\partial^2G_s/\partial x^2(0)=2s$,
the critical equation in $x$ has a unique
analytic solution $x=\xi(y,z)$ with
$\xi$ divisible by $yz$. The parametric
holomorphic Morse change in $x$ writes
\[
 G_s=X^2+y^2H(y,z).
\]
Divisibility follows directly on substituting
$\xi$: each of $\xi yz,\xi^p,s\xi^2$
is divisible by $y^2$, as are the two
pure $y$ terms. Moreover
\[
 H_y(0,0)=
 \begin{cases}1+s&q=3,\\s&q>3,\end{cases}
 \qquad H(0,z)=-z^2/(4s).
\]
The finite
birational map
\[
 \{v^2+H(y,z)=0\}\longrightarrow\{G_s=0\},
 \qquad (v,y,z)\longmapsto(X=yv,y,z),
\]
followed by the inverse Morse coordinate
change, is the normalization. Its source
is smooth because $H_y(0,0)\ne0$,
and the map is finite because $v$
satisfies the displayed monic equation.
It is birational on $y\ne0$ and is
therefore the normalization by normality
of its source. At the origin its
differential has rank one. The normalization of a normal-crossing
germ consists of coordinate-plane inclusions, all of differential rank
two. The normalization isomorphism induced by an ambient biholomorphism
preserves these ranks, so this rank-one map excludes normal crossings. The local construction in Proposition~\ref{night:P4.20} produces
a spherical class of nonzero square, and its boundary transport gives
the contradiction. The transverse Hessians and the restrictions
$H(0,z)$ and $H_y(0,0)$ follow by collecting the degree-two terms
and substituting the critical equation
$yz+p\xi^{p-1}+2s\xi=0$.

\item The remaining germs in Table~1 are smooth, ordinary double,
and triple normal crossing. The linear coordinates
$(u,v)=(x+iy,x-iy)$ give $uv=0$ for the double germ;
the triple germ is $xyz=0$. Their normalization morphisms are the
inclusions of their disjoint smooth coordinate-plane branches.
Proposition~\ref{night:P4.6} applies to each compact support so obtained.
When every reduced support is semi-log-canonical, the same local
conclusion at every fibre supplies the hypotheses of
Propositions~\ref{night:P4.7}--\ref{night:P4.8}.

\end{enumerate}
\end{proof}